% ECONOMICS_PROBLEM: {"id": "D71-652", "title": "Core nonemptiness: approval counts versus general additive scores", "jel_code": "D71", "source_id": "OP-0196"}
% !TeX program = xelatex
\documentclass[11pt,a4paper]{article}
\usepackage[margin=23mm]{geometry}
\usepackage{fontspec}
\setmainfont{Latin Modern Roman}
\usepackage{amsmath,amssymb,mathtools}
\usepackage{xcolor,enumitem,booktabs,longtable,array,xurl,hyperref}
\definecolor{muted}{HTML}{53636C}
\hypersetup{colorlinks=true,urlcolor=blue,linkcolor=blue}
\setlength{\parindent}{0pt}
\setlength{\parskip}{5pt}
\newcommand{\R}{\mathbb R}
\newcommand{\E}{\mathbb E}
\newcommand{\N}{\mathbb N}
\newcommand{\ind}{\mathbf 1}
\newcommand{\Var}{\operatorname{Var}}
\newcommand{\Cov}{\operatorname{Cov}}
\newcommand{\supp}{\operatorname{supp}}
\DeclareMathOperator*{\argmin}{arg\,min}
\DeclareMathOperator*{\argmax}{arg\,max}
\newcommand{\entrymeta}[3]{{\sffamily\small\color{muted}\textbf{\texttt{#1}}\quad|\quad #2\par #3\par}}
\newcommand{\Problemref}[1]{\texttt{#1}}
\newcommand{\JELref}[2]{\texttt{#2} (JEL \texttt{#1})}
\begin{document}
\section*{Core nonemptiness: approval counts versus general additive scores}
\textbf{Problem ID:} \texttt{D71-652}\quad \textbf{JEL:} \texttt{D71}\par
% BEGIN ECONOMICS STATEMENT
\label{problem:OP-0196}
% GAME_THEORY_PROBLEM: {"local_id": "GT-066", "original_number": 66, "kind": "H", "entry_kind": "statement_only", "published": false, "review_required": true, "category": "game_theory"}
\label{game:GT-066}
{\sffamily\small\color{muted}
\noindent\textbf{ID: GT-066.}\quad\textbf{Category:} Committee voting.
\quad\textbf{Kind: H.}\quad\textbf{Status:} General-score strengthening false;
approval-count interpretation duplicates Entry~\ref{game:GT-072}.
\par}
\subsection*{Definitions and mathematical statement}
Let \(N=\{1,\ldots,n\}\), \(n\geq 1\), let \(C\) be a finite candidate
set, and let \(0\leq k\leq |C|\) be an integer.
Give voter \(i\) nonnegative rational scores \(a_{ic}\), and set
\(u_i(T)=\sum_{c\in T}a_{ic}\).
A committee \(W\subseteq C\), \(|W|=k\), is core-stable if there are no
nonempty \(S\subseteq N\) and \(T\subseteq C\) with
\[
 |T|\leq \frac{|S|}{n}k,\qquad
 u_i(T)>u_i(W)\quad\text{for every }i\in S.
\]

\paragraph{Historical claim, now refuted.}
The assertion that every such general-score instance has a core-stable
committee is false. For example, take \(n=6\), \(k=3\),
\(C=\{c_1,\ldots,c_6\}\), and the score matrix
\[
 (a_{ic})=
 \begin{pmatrix}
 2&1&0&0&0&0\\
 0&2&1&0&0&0\\
 1&0&2&0&0&0\\
 0&0&0&2&1&0\\
 0&0&0&0&2&1\\
 0&0&0&1&0&2
 \end{pmatrix}.
\]
Every three-member committee uses at most one candidate from one of the
two triples. Relabel that triple cyclically so that its second and third
candidates are absent. The voters corresponding to the second and third
rows of that triple block the committee by choosing its third candidate:
each strictly improves, and one seat equals their proportional budget
\((2/6)3\).

\subsection*{Short English statement}
General additive candidate scores can produce an empty core.
The approval-count nonemptiness question has a positive resolution reported
in the September 2026 preprint cited in Entry~\ref{game:GT-072}.

\subsection*{Variants and scope boundaries}
For an approval set \(B_i\subseteq C\), approval-count utility is
\(u_i(T)=|B_i\cap T|\), equivalently \(a_{ic}\in\{0,1\}\).
It is additive, and is not the unit-demand utility
\(\mathbf 1_{\{B_i\cap T\ne\varnothing\}}\).
Allowing arbitrary positive scores changes the problem and invalidates
universal nonemptiness.

\subsection*{Source relationship and status}
The catalogue conflated these models. [S1, Section 7.2] distinguishes
them; [S2, Example 2] supplies the counterexample instantiated above.
The original entry is retained as a corrected historical claim.

\subsection*{References}
\begingroup\footnotesize\raggedright
\noindent[S1] Patrick Becker, Matthias Greger, and Dominik Peters,
\emph{Core Existence in Approval-Based Committee Elections with up to
Seven Voter Types}, 2026, arXiv:2605.06194v2.
\url{https://arxiv.org/html/2605.06194v2}.
Locators: Section 2.1, Definition 2.1; Section 7.2, \emph{Additive Valuations}.

\noindent[S2] Dominik Peters, Grzegorz Pierczy\'nski, and Piotr Skowron,
\emph{Proportional Participatory Budgeting with Additive Utilities},
NeurIPS 2021; arXiv:2008.13276v2, 2022.
\url{https://arxiv.org/html/2008.13276v2}.
Locator: Section 3.2, Example 2.
\par\endgroup

\subsection*{Common conventions: Source mathematical conventions}

\textbf{Finite games.} For a finite set $A$, $\Delta(A)$ is its probability
simplex. Players choose independent mixed strategies in Nash equilibrium;
correlation is allowed only when explicitly part of the solution concept.
In a game with payoff functions $u_i$ and strategy profile $x$, an additive
$\varepsilon$ Nash equilibrium satisfies
\[
 u_i(x)\geq u_i(x_i',x_{-i})-\varepsilon
 \quad\text{for every player $i$ and every admissible deviation $x_i'$}.
\]
For numerical approximation, payoffs are normalized as locally specified.
An $\varepsilon$ payoff residual is distinct from distance at most
$\varepsilon$ to the set of exact equilibria. Point convergence, convergence
of empirical play, and convergence in distance to an equilibrium set are
also distinct.

\textbf{Computational representations.} Unless stated otherwise, a complexity
question takes rational data with binary encoding; polynomial time is measured
in total bit length. A fixed-accuracy polynomial algorithm is not necessarily
polynomial in $1/\varepsilon$, and neither is necessarily polynomial in
$\log(1/\varepsilon)$. Succinct games and value, demand, or sampling oracles
must declare their interfaces. Exact real solutions require an explicit
representation; an unspecified list of arbitrary real numbers is not a
finite computational output. A claimed algorithmic barrier is meaningful
only for its specified model and reductions.

\textbf{Dynamic games.} Histories, information, observation, and allowable
randomization are part of the model. A stationary strategy depends only on
the current state; a positional strategy in a deterministic graph game
chooses one action at each controlled vertex. Nash equilibrium at an initial
state and equilibrium after every history are different requirements.
An expectation of a pathwise $\liminf$ payoff, a $\liminf$ of expected averages,
a limiting expected average, and a uniform finite-horizon equilibrium are
different criteria. None is silently substituted for another.

\textbf{Allocation and mechanisms.} A complete allocation assigns every item
exactly once unless otherwise stated. Zero-valued goods, negative values,
capacity constraints, feasibility, lotteries, and copies of goods affect
fairness definitions and are specified locally. Dominant-strategy incentive
compatibility, Bayesian incentive compatibility, universal truthfulness,
and truthfulness in expectation impose different inequalities. Whenever
an objective ratio has zero denominator, its convention is stated or those
instances are excluded.

\textbf{Infinite-dimensional models.} Expectations, integrals, stochastic
equations, and optimization objectives require their local regularity,
integrability, filtrations, and admissible domains. A backward--forward
mean-field system is a boundary-value problem; it is not an initial-value
semiflow without an additional construction. Long-time, large-population,
and vanishing-noise limits require an order of limits and a convergence
topology. If a minimum need not be attained, the objective is its infimum
or supremum and the approximation requirement is stated separately.
% END ECONOMICS STATEMENT
\end{document}
